\section{Introduction}\label{sec:introduction}

Can the free Ising state on a regular tree be generated from independent
vertex labels without choosing a root? The broadcast description of the
state begins at a vertex and propagates a random spin along the edges.
Its law does not depend on that vertex. A factor construction must also
remove the distinguished vertex from the rule that generates the spins.

Let $T_d=(\mathcal V,\mathcal E)$ be the infinite $d$-regular tree,
$d\ge3$, and put $b=d-1$ and $\theta=\tanh\beta$ for $\beta\ge0$.
The free zero-field Ising law $\mu_{d,\beta}$ is obtained by choosing a
fair spin at one vertex and, independently along edges directed away
from it, letting a child agree with its parent with probability
$(1+\theta)/2$. Its marginal on any finite connected subtree $D$ is
\[
 \mu_{d,\beta}(\sigma|_D=\eta)
   =\frac12\prod_{\{u,v\}\in\mathcal E(D)}
                      \frac{1+\theta\eta_u\eta_v}{2}.
\]
This expression also verifies that the law is invariant under all tree
automorphisms and independent of the chosen broadcast root.

We call a spin law a \emph{factor of IID} if it is the law of
$\Phi(U)$ for independent uniform labels $(U_v)_{v\in\mathcal V}$ and
a measurable map $\Phi:[0,1]^{\mathcal V}\to\{-1,1\}^{\mathcal V}$
satisfying
\begin{equation}\label{eq:equivariance}
 \Phi(gU)=g\Phi(U)\quad\text{almost surely for each fixed }
                 g\in\Aut(T_d),\qquad (gx)_v=x_{g^{-1}v}.
\end{equation}
The exceptional null set may depend on $g$. For spin laws this convention
is equivalent to requiring a Borel factor equivariant on every input:
Lemma~\ref{lem:pointwise-spin-factor} constructs such a version.
The input contains no distinguished root or end. Measurability does not require a spin to be
determined by finitely many labels; no finitary coding or coding-radius
estimate is asserted.

\begin{theorem}\label{thm:main}\label{thm:sharp}
For every integer $d\ge3$ and every $\beta\ge0$, the free zero-field
Ising law $\mu_{d,\beta}$ is a factor of IID in the sense of
\eqref{eq:equivariance} if and only if
\[
                         \tanh\beta\le(d-1)^{-1/2}.
\]
In this range the factor can be chosen to commute with every
automorphism on every label input.
\end{theorem}

The construction proves the positive implication, including equality.
The strict converse is known: the obstruction beyond the reconstruction
threshold is due to Sly, as recorded and proved in
Lyons~\cite[Section~3 and Corollary~3.2]{Lyons}. We use the quantitative correlation bound
of Backhausz, Szegedy and Vir\'ag~\cite[Theorem~3.1 in
arXiv:1305.6784v1]{BSV} to give that implication below.
The map is obtained for each fixed $d$ and $\beta$ in the permitted
range; no parameter-uniform or computational bound is asserted.

For the free state on $T_d$, reconstruction from distant levels is
possible exactly when $b\theta^2>1$; equivalently, the free Gibbs state
is extremal exactly when $b\theta^2\le1$. For the regular-tree result,
including equality, see the Bethe-lattice work of Bleher, Ruiz, and
Zagrebnov~\cite{BleherRuizZagrebnov1995} and
Ioffe~\cite{Ioffe1996Bethe}. Ioffe~\cite{Ioffe1996GeneralTrees} and
Evans, Kenyon, Peres, and Schulman~\cite{EvansKenyonPeresSchulman2000}
developed branching-number criteria on general trees away from equality.
Pemantle and Peres~\cite[Theorem~2.1]{PemantlePeres} later gave an
electrical-capacity criterion for free-state extremality for couplings
bounded above and away from zero, which also covers critical cases.

Reconstruction asks whether distant spins retain information about a
fixed spin. The factor question asks for a rule that generates the
entire law from independent labels. Nonreconstruction supplies no such
rule, and a weak limit of factor laws need not itself be a factor
\cite[Corollary~3.3]{Lyons}. Thus equality requires a construction that retains
measurable dependence on its independent input.

In the uniqueness range $\theta\le b^{-1}$, including equality, a
finite-cluster construction gives the factor; see
Lyons~\cite[Section~3]{Lyons}. Nam, Sly and
Zhang~\cite[Theorem~1 and Section~2.2 in arXiv:2012.09484v2]{NSZ}
used posterior-drift stochastic differential equations driven by IID
Brownian motions to construct the free Ising law as a factor for
$d\ge d_0$ and $\theta\le c/\sqrt{d-1}$, for unspecified absolute
constants $c,d_0>0$. Their construction takes root-independent limits
of finite-volume solutions.
Theorem~\ref{thm:main} proves their conjecture for every $d\ge3$,
including equality and their pointwise convention for equivariance.

\paragraph{The known converse.}
For a centered variance-one real vertex factor, the BSV bound is
\[
 |\E[Z_oZ_v]|\le
       \left(1+\frac{d-2}{d}n\right)b^{-n/2},
                  \qquad \dist(o,v)=n.
\]
It applies also under~\eqref{eq:equivariance}. To check this,
condition the observable at $o$ on the labels in its radius-$R$ ball
and average a Borel version over the finite automorphism group of that
rooted ball. Every such permutation extends to a tree automorphism
fixing $o$, so uniqueness of conditional expectation leaves this
average unchanged almost surely. It now defines an exactly invariant
local rule. Apply the same rule at every vertex. These local factors
converge in $L^2$ at each vertex to the original factor, because the
balls exhaust the label coordinates. For all sufficiently large $R$
their variances are positive. Center and normalize them, apply BSV,
and pass to the limit in variances and covariances on this same IID
probability space. Since the Ising spins have mean zero, variance one
and correlation $\theta^n$, the bound would require
$(\theta\sqrt b)^n\le1+(d-2)n/d$. This is impossible when
$\theta>b^{-1/2}$.

\subsection{Why independent innovations determine the spins}

We follow the observation framework of Nam--Sly--Zhang. On an auxiliary
probability space, sample the desired spin configuration $\sigma$ and
independent standard Brownian motions $B_v$, and observe
\[
 X_v(t)=t\sigma_v+B_v(t),\qquad
 m_v(t)=\E[\sigma_v\mid X_u(s):u\in\mathcal V,\ 0\le s\le t].
\]
The finite-dimensional posterior evolution also belongs to stochastic
localization. Eldan~\cite{Eldan2013} developed the underlying localization
scheme, and El Alaoui and Montanari~\cite[Section~3, Theorem~2]{ElAlaouiMontanari2022}
give the fixed-covariance Gaussian-observation formulation used here.
Subtracting the conditional drift gives
\[
                      W_v(t)=X_v(t)-\int_0^t m_v(s)\,ds.
\]
These innovations are independent Brownian motions in the observation
filtration. This property alone does not make $X$ a function of $W$:
the original construction still used the spins. Classical innovations
theory separates the Brownian property from generation of the
observation filtration~\cite[Lemma~2.2 and Remark~2.2]{FujisakiKallianpurKunita1972}.
Proving that the innovations determine $X$ is the central step here.

A companion construction for the free uniform spanning forest uses a
uniform Lipschitz response bound in the coordinate supremum norm to
invert the innovations causally for every continuous input
\cite[Section~4]{OpenAIFUSF2026}.
Here recovery instead follows from equality of conditional copies under
the observation law.

Take two observations $X,X'$ independently conditional on the entire
innovation array $W$. They have the same marginal law and share the
same innovations. Conditioning on the full paths could in principle
expose their future increments through an observation past. We prove
that it does not: the conditional law of either observation past given
all of $W$ depends only on the past of $W$. Hence $W$ remains Brownian
in the joined filtration, and the equations for both observations use
the same driving Brownian motions. This is the compatibility principle
for conditional copies~\cite[Lemmas~2.11--2.12 in arXiv:1305.6747v2]{Kurtz2014};
we prove the required statement
and the transport of the stochastic integrals directly.

We then show that these particular copies agree. For each coordinate,
conditional independence gives
\[
 \E[(X_v(t)-X'_v(t))^2\mid W]
                  =2\Var(X_v(t)\mid W).
\]
Equality therefore forces the conditional law of $X$ given
$W$ to be a point mass. Finally
\[
                         \frac{X_v(n)}n\longrightarrow\sigma_v
                           \quad\text{almost surely}
\]
recovers the spins from $W$. A common measurable sampler converts one
uniform label at each vertex into its Brownian path. This gives a spin
factor satisfying~\eqref{eq:equivariance}; finite-ball averaging then
selects an everywhere equivariant version with the same law.

\subsection{What changes at the critical point}

The tree gives a concrete equation for comparing the copies. Set
$h_v=\operatorname{atanh}m_v$. If $T_{j\to i}$ is the component containing
$j$ after deleting $ij$, define its field by
\[
 h_{j\to i}(t)=\operatorname{atanh}
   \E[\sigma_j\mid X_u(s):u\in T_{j\to i},\ 0\le s\le t],
\]
using the unconditional prior on that component.
For a path $v=v_0,\ldots,v_r=u$, the field-response
coefficient is
\[
 c_{vu}=\theta^r\prod_{i=1}^r a(h_{v_i\to v_{i-1}}),
 \qquad
 a(z)=\frac{1-\tanh^2z}{1-\theta^2\tanh^2z},\qquad c_{vv}=1.
\]
The innovation equation is
\[
 dh_v=\sum_u c_{vu}\,dW_u+m_v A_v\,dt,
                  \qquad A_v=\sum_u c_{vu}^2.
\]
Its infinite sum requires proof. At criticality we obtain
$\E A_v(t)\le C/\min(\sqrt t,1)$ for $t>0$, so the row is integrable
down to time zero. An exact covariance-energy identity and a testing
argument establish the equation without a residual martingale.

For the conditional copies, let
$V(t)=\E[(h_o(t)-h'_o(t))^2]$; invariance makes the choice of $o$
irrelevant. It starts at zero. Below the threshold, the squared factor
$\theta^{2r}$ and the shell size $d b^{r-1}$ give a summable geometric
series, yielding $V'\le CV$ and thus equality of the copies.
At $b\theta^2=1$ those two factors cancel exactly.

Three estimates replace the lost geometric decay. First, an intact side
branch attenuates a path product on the length scale $t^{-1/2}$ at small
times. Second, spatial cancellation bounds the variance of a sum of
differences between the two copies' $\log a(h_{j\to i})$ factors by
the path length times the variance at one edge. The spectral
bound is proved for differences $f(X)-f(X')$, with the same
square-integrable equivariant vertex observable $f$ applied to each copy, and their closed
linear span. No spectral bound for every function of the pair is assumed.
The operator argument is related to the correlation methods of
Backhausz--Szegedy--Vir\'ag~\cite{BSV} and
Backhausz--Gerencs\'er--Harangi--Vizer~\cite{BGHV}.
Third, concentration of path averages makes the attenuation
usable with discrepancy weights that depend on the same observations.
Both path orientations are needed in this averaging argument.

These estimates leave one small-time difficulty: the leading noise
term is of order $S(t)/t$, where
\[
 S(t)=\E\Bigl[\bigl(\log a(h_{j\to i}(t))
                    -\log a(h'_{j\to i}(t))\bigr)^2\Bigr].
\]
The elementary bound $S\le CV$ would leave a nonintegrable coefficient.
We prove instead a small-field tail estimate and use the vanishing
derivative of $\log a$ at zero to improve that comparison. The complete
noise and drift calculation then gives an integrable Osgood-type
inequality for $V$, forcing $V=0$. The moment estimate is proved beside
this final comparison, where its purpose becomes explicit.

\paragraph{Notation and organization.}
The case $\theta=0$ is already IID, so the proof below assumes
$0<\theta\le b^{-1/2}$. Constants $c,C>0$ may depend on the fixed
$d,\theta$ and may change between occurrences; a subscript $T$ permits
additional dependence on a fixed finite time horizon. Every path is
simple unless specified otherwise, and all logarithms are natural.
Roots used to describe subtrees or compute norms are auxiliary choices
in the proof, not inputs to the factor.

Sections~\ref{sec:observations}--\ref{sec:innovations} develop the
posterior identities, critical information bounds, and infinite
innovation equation. Section~\ref{sec:coupling} constructs the conditional
copies and proves their equality below the threshold.
Sections~\ref{sec:spectrum} and~\ref{sec:path-averages} establish spatial
cancellation and concentration of path averages.
Section~\ref{sec:uniqueness-factor} proves the small-field tail estimate,
completes the critical comparison, and recovers the factor from the
innovations.
